\documentclass[12pt]{article}
\usepackage[utf8]{inputenc}
\usepackage[T1]{fontenc}
\usepackage[margin=1in]{geometry}
\usepackage{setspace}
\usepackage{xcolor}
\usepackage[colorlinks=true, urlcolor=blue!60!black, linkcolor=blue!60!black]{hyperref}
\setlength{\parindent}{2em}
\setlength{\parskip}{1.25em}
\renewcommand{\baselinestretch}{1.0}
\begin{document}
\begin{center}
{\LARGE \textit{United States v. Hemani}, No. 24-1234 (2026)\par}
\vspace{0.8em}
{\large Case Brief --- Four-Page Edition\par}
\end{center}
\vspace{1em}

\section{Memory Jogger}

A unanimous Supreme Court held that prosecuting a regular marijuana user under 18 U.S.C. \S~922(g)(3) solely for possessing a firearm at home violates the Second Amendment, because historical ``habitual drunkard'' laws targeted incapacitated persons for different reasons and only after pre-deprivation process (\href{https://www.supremecourt.gov/opinions/25pdf/24-1234_g2bh.pdf}{Slip Op.}).

\section{Detailed Case Facts}

In August 2022, federal agents searched the Dallas-area home of Ali Danial Hemani, a U.S.--Pakistani dual citizen, on suspicion of terrorism-related activity. Hemani cooperated: he surrendered a Glock pistol, pointed agents to marijuana, and admitted using it ``about every other day.'' A grand jury indicted him on a single count under \S~922(g)(3), which bars any ``unlawful user of or addicted to any controlled substance'' from possessing a firearm (up to 15 years' imprisonment). The government brought no terrorism, cocaine, or trafficking charges and did not allege that Hemani was addicted, intoxicated while armed, or dangerous.

\section{Procedural History}

The Eastern District of Texas dismissed the indictment on Second Amendment grounds. The Fifth Circuit affirmed, applying its precedent in \textit{United States v. Connelly}, 117 F.4th 269 (5th Cir. 2024), which barred applying \S~922(g)(3) to a sober regular marijuana user. The Supreme Court granted the government's petition, heard argument on March 2, 2026, and decided the case on June 18, 2026.

\section{Judicial Votes}

Affirmed, \textbf{9--0}.
\begin{itemize}
\item \textbf{Majority (7):} Justice Gorsuch, joined by Chief Justice Roberts and Justices Thomas, Sotomayor, Kavanaugh, Barrett, and Jackson.
\item \textbf{Concurrences:} Justice Thomas (Commerce Clause doubts about \S~922(g)); Justice Jackson, joined by Justice Sotomayor (criticizing \textit{Bruen}).
\item \textbf{Concurrence in the judgment:} Justice Alito, joined by Justice Kagan.
\end{itemize}

\section{Holding}

Applying \textit{Bruen} and \textit{Rahimi}, the government may not strip an individual of his Second Amendment right under \S~922(g)(3)'s ``unlawful user'' provision solely because he regularly uses a controlled substance such as marijuana, without pre-deprivation process or proof of incapacitation or dangerousness. Historical habitual-drunkard laws are not ``relevantly similar'' in their ``why'' or their ``how.'' The Court expressly reserved bans on addicts or the presently intoxicated, targeted laws for drugs proven to create special risks, \S\S~922(g)(1) and (g)(4), and prosecutions supported by individualized proof of dangerousness.

\section{Analysis of Opinions}

\textbf{Majority (Gorsuch).} Because \S~922(g)(3) burdens conduct the Second Amendment's text covers, the government had to show its prosecution fit the Nation's historical tradition. Its analogues---vagrancy, civil-commitment, and surety-of-good-behavior laws for ``habitual drunkards''---failed on every metric. They targeted different people: in an era of copious drinking, a habitual drunkard was someone so persistently intoxicated as to be incapable of managing his affairs, not a regular drinker, whereas the government's reading would reach a student using a friend's Adderall. They served different purposes (labor norms, protecting families, policing moral scandal). And they operated differently, requiring judicial process before any deprivation. Finally, the government's dangerousness rationale was undercut by DOJ non-enforcement, widespread state legalization, and the 2026 rescheduling of some marijuana products.

\textbf{Thomas (concurring).} Mere intrastate gun possession is non-economic activity, so \S~922(g) may exceed Congress's commerce power under \textit{Lopez} and \textit{Morrison}.

\textbf{Jackson, with Sotomayor (concurring).} The Court should retire ``the failed \textit{Bruen} experiment'' and return to means-end scrutiny.

\textbf{Alito, with Kagan (concurring in the judgment).} Would affirm on one narrow ground: the record showed nothing about the quantity, potency, or impairing effect of Hemani's marijuana use.

\section{Examples: Future Applications}

\begin{enumerate}
\item \textbf{Same-Side (Unconstitutional): Medical-Marijuana Homeowner.} A Colorado resident who uses state-legal medical marijuana twice a week keeps a locked shotgun at home and has no record of violence. A \S~922(g)(3) prosecution resting solely on her regular use fails, absent process or individualized proof of dangerousness.
\item \textbf{Same-Side (Unconstitutional): Off-Label Adderall User.} A graduate student who takes a friend's unprescribed Adderall to study and keeps a handgun at home, with no sign of erratic behavior---the very overbreadth example Justice Gorsuch used.
\item \textbf{Opposite-Side (Likely Constitutional): Armed While Intoxicated.} A chronic user, actively high and carrying a loaded handgun, threatens bystanders during a street dispute. \textit{Hemani} reserved bans on possession while intoxicated and prosecutions backed by individualized proof of dangerousness.
\item \textbf{Opposite-Side (Likely Constitutional): Judicial Commitment of an Addict.} After an adversarial hearing, a probate judge finds a person severely addicted to fentanyl and suspends his firearm rights during rehabilitation. The measure targets genuine incapacitation and provides pre-deprivation process, mirroring Founding-era commitment laws.
\item \textbf{Fence-Sitter (Close Case): Frequent Methamphetamine User.} A defendant who uses methamphetamine several times a week flees police and is arrested with a pistol, but the government offers no evidence that his use made him violent or incapacitated. Methamphetamine is far more dangerous than marijuana, yet mere frequency of use is not enough on its own.
\end{enumerate}

\section{Critique}

\textbf{Originalist and textualist views.} Some originalists praised \textit{Hemani} for refusing to let the government strip a core right on its bare ``say-so'': history fixes the permissible reasons for disarmament, and an as-applied record checks pretext. Others faulted the Court for ducking the vagueness of ``unlawful user'' and leaving lower courts little portable guidance; Justice Thomas's concurrence points to a deeper Commerce Clause problem.

\textbf{Progressive and doctrinal views.} Echoing Justice Jackson, critics argue that \textit{Hemani} exposes \textit{Bruen}'s incoherence: the majority engaged in covert means-end reasoning, and its narrow holding invites case-by-case fights over drug potency and individualized dangerousness. The case pits individual liberty and notice against administrable, bright-line public-safety rules.

\section{Key Quotations}

\begin{quote}
``And the habitual drunkard laws on which the government relies here differ dramatically from \S~922(g)(3)'s unlawful user provision on every single metric the government invites us to consider: They targeted different kinds of people, did so for different purposes, and operated in different ways.''
\end{quote}
\noindent --- Justice Neil M. Gorsuch, Majority Opinion, slip op. at 7.

\begin{quote}
``All of which leaves it awkwardly positioned to suggest that the millions of Americans who now regularly use marijuana are categorically and unusually dangerous.''
\end{quote}
\noindent --- Justice Neil M. Gorsuch, Majority Opinion, slip op. at 17.

\begin{quote}
``In a future case that squarely presents the question, we should consider whether to retire the failed Bruen experiment and return to an explicit assessment of Congress's ends and means when deciding the constitutionality of firearm restrictions.''
\end{quote}
\noindent --- Justice Ketanji Brown Jackson, Concurring Opinion, slip op. at 3.

\end{document}
